\section{Discussion} \label{discussion}

Evolution depends on the joint action of natural selection and the generation of heritable variation, yet the two are usually studied with local tools (the selection gradient and the G- or M-matrix) that summarize only the neighborhood of a population's current state. On rugged, multi-peaked landscapes, and with nonlinear genotype-phenotype maps where the M-matrix evolves as the population mean shifts, these local quantities need not anticipate where evolution ends up. We here investigated combinations of adaptive landscapes and genotype-phenotype maps in order of ascending complexity. Our three main findings are 1) that developmental bias can distort the domains of attraction of alternative evolutionary optima, including in nonlinear genotype-phenotype maps where developmental bias varies over phenotype space, 2) that the influence of developmental bias depends on the complexity of the adaptive landscape, and 3) that in nonlinear genotype-phenotype maps developmental bias can evolve, leading to divergent evolutionary outcomes even if starting conditions are identical.  

We first demonstrated that even linear genotype-phenotype maps, which induce constant M-matrices over phenotype space, can reshape the domains of attraction of alternative evolutionary optima. Developmental bias therefore not only leads to a temporary distortion of evolutionary trajectories, but also affects long-term evolutionary outcomes. The same occurs when phenotypes are under the control of gene regulatory networks (GRN), which induce nonlinear genotype-phenotype maps and M-matrices that vary over phenotype space. The GRNs are modeled as a minimal dynamical system, and thus as a mechanistic representation of development \parencite{jaeger2014bioattractors, milocco2024utilizing}, departing from standard quantitative-genetic treatments, where phenotypes are direct (often additive) functions of underlying parameters (\cite{Arnold2008-pc, jones2007mutation,Jones2014-wj,jiang2025alternative}, but see \cite{gjuvsland2007statistical, hether2014genetic}). The distortion to the domains of attraction varies over phenotype space as a function of the nonlinear genotype-phenotype map, and is the strongest where M-matrices become very eccentric. 
 
 Previous work in quantitative genetics observed that standing genetic variation can drive evolutionary outcomes \parencite{steppan2002comparative, melo2016modularity}, without however quantifying how domains of attraction are altered. The only other authors operating within the framework of quantitative genetics that traced out a separatrix based on genetic correlations are \cite{price1993peak}. However, they assumed genetic correlations to be fixed rather than connecting them to the consistent effect of developmental bias induced by the genotype-phenotype map. Characterizations of the genotype-phenotype map may have greater potential for longer-term evolutionary prediction than the G-matrix, which combines mutation and recombination with selection and demographic stochasticity, and may as a result be more variable (e.g. \cite{bjorklund2013rapid}). Conversely, the more common empirical observation of relative G-matrix stability---which helps to explain the often surprising effectiveness of macroevolutionary predictions based on standing variation \parencite{Tsuboi2024-gm}---suggests that this variation is consistently anchored by underlying developmental mechanisms. This underscores the opportunity for quantitative genetics to integrate developmental mechanisms into its framework \parencite{hansen2008macroevolutionary, milocco2026bridging}.  The connection between  the M-matrix and the Jacobian of the dynamical system (equation \ref{eqn:MandJ}) provides one possible avenue to do this. Indeed, structural properties of the Jacobian can constrain how $\mathbf{M}$ evolves, and these constraints could be used to model the evolution of the M-matrix. For example, it has been recently suggested that evolved developmental systems exhibit Jacobians with dominant modes \parencite{russo2025soft, milocco2026bridging, furusawa2018formation}, where a small number of eigenvectors govern large-scale behavior. On the other hand, we are quick to note that genotype-phenotype maps will generally not be bijective one-to-one mappings, but rather many-to-one, greatly complicating the analysis. At the same time, mechanistic insights on development may provide opportunities for simplification \parencite{gonzalez2024mathematical}. 

Our second result is that the influence of developmental bias on evolutionary outcomes depends on the complexity of the adaptive landscape. We varied landscape complexity defined as the frequency of the Perlin noise, resulting in multipeaked landscapes, and measured whether populations converged on the peak predicted by selection alone or diverged toward alternative optima. In single-peaked landscapes, selection determines the outcome: all populations reach the same optimum, and developmental bias only affects their trajectory, as expected (\FIG[A]{fig1}). As landscape complexity increases, however, the optimum reached is increasingly determined by developmental bias (\FIG[]{fig4}). Furthermore, with increasing landscape complexity, the alignment between $M_{max}$, the dominant eigenvector of $\mathbf{M}$, along which most mutational and genetic variation occurs, and the direction of evolutionary divergence becomes increasingly well-aligned (\FIG[]{fig5}). Such alignment between variation and long-term divergence has been reported in several system \parencite{Houle2017-ph, rohner2023developmental, holstad2024evolvability}, but the causes of this alignment are not fully resolved, and the extent to which this generalizes to other organs and organisms is currently unclear \parencite{rohner2025developmental}. Our findings suggest that the interplay between the statistical features of the adaptive landscape and the structure of underlying developmental biases could provide an explanation for these ambivalent empirical observations.

Our third result is that, as a result of the evolution of developmental bias discussed above, its long-term consequences cannot always be predicted from a bias measured at a single generation. In our simulations, $\mathbf{M}$  emerges from GRN dynamics rather than being imposed, and therefore changes over the course of evolution. Differences in the way $\mathbf{M}$ evolves are sufficient to drive divergence: populations starting with very similar M-matrices and evolving on the same adaptive landscape can reach different optima (\FIG{fig6}). While the breeder’s equation establishes that standing variation biases short-term responses \parencite{lande1979quantitative}, we show that evolving developmental bias can determine which adaptive peak is ultimately reached and, therefore, dictate long-term evolution.

Our model of development does not include hard developmental constraints that render some phenotypes entirely unreachable. Indeed, the gene regulatory networks we use can generate phenotypes across all positive values of gene expression. Divergence here therefore arises only from differences in the relative accessibility of phenotypic directions, rather than the exclusion of any direction. In real systems, nonlinear developmental interactions can exclude regions of phenotype space, such as inaccessible tooth \parencite{salazar2010computational} or beak \parencite{al2021geometry} shapes, and such constraints would be expected to amplify the effects of development on evolutionary divergence \parencite{milocco2020evolution, rice2002general, gonzalez2024mathematical}. The magnitude of divergence we measure should therefore be read as a conservative estimate, since real developmental systems combine bias of the kind modeled here with hard constraints leading to the exclusion of regions of phenotype space.

In summary, we find that the balance between selection and variation as difference makers in evolution is context-dependent, with developmental bias exerting increasing influence under complex landscapes. Understanding this interplay between adaptive landscape and developmental bias has direct implications for evolutionary prediction: if developmental bias is negligible, characterizing the adaptive landscape is enough for reliable prediction; if it is not, then the structure and evolution of developmental bias, the systematic tendency of developmental systems to produce some variants more readily than others, becomes central. In such situations in particular, a mechanistic, dynamical understanding of variation is essential, and not just supplementary, to explaining and predicting evolutionary outcomes.